\begin{question}{2.7}

Cho các vector $X_1, X_2, X_3, X_4$ độc lập và có cùng phân phối
\[
N_p(\mu,\Sigma).
\]

\begin{enumerate}
\item[(a)] Tìm các phân phối lề của các vector ngẫu nhiên
\[
V_1=\frac{1}{4}X_1-\frac{1}{4}X_2+\frac{1}{4}X_3-\frac{1}{4}X_4
\]

và
\[
V_2=\frac{1}{4}X_1+\frac{1}{4}X_2-\frac{1}{4}X_3-\frac{1}{4}X_4.
\]
\item[(b)] Tìm hàm phân phối đồng thời của $V_1$ và $V_2$.
\end{enumerate}

\textbf{Lời giải:}

Cho
\[
X_1, X_2, X_3, X_4 \sim N_p(\mu,\Sigma)
\]
độc lập với nhau.

Đặt
\[
V_1=\frac{1}{4}X_1-\frac{1}{4}X_2+\frac{1}{4}X_3-\frac{1}{4}X_4
\]

\[
V_2=\frac{1}{4}X_1+\frac{1}{4}X_2-\frac{1}{4}X_3-\frac{1}{4}X_4
\]

%--------2.7a--------
\begin{subquestion}
\begin{answer}
\textbf{Phân phối lề của $V_1$ và $V_2$}

Do tổ hợp tuyến tính của các vector chuẩn nhiều chiều vẫn là chuẩn nhiều chiều nên $V_1, V_2$ đều có phân phối chuẩn nhiều chiều.

Tính kỳ vọng:
\[
E(V_1)=\frac{1}{4}\mu-\frac{1}{4}\mu+\frac{1}{4}\mu-\frac{1}{4}\mu=0
\]

\[
E(V_2)=\frac{1}{4}\mu+\frac{1}{4}\mu-\frac{1}{4}\mu-\frac{1}{4}\mu=0
\]

Tính hiệp phương sai:

Vì các $X_i$ độc lập nên
\[
\mathrm{Var}(aX)=a^2\Sigma
\]

Do đó
\[
\mathrm{Var}(V_1)=\left(\frac{1}{4}\right)^2(\Sigma+\Sigma+\Sigma+\Sigma)
=\frac{1}{4}\Sigma
\]

\[
\mathrm{Var}(V_2)=\left(\frac{1}{4}\right)^2(\Sigma+\Sigma+\Sigma+\Sigma)
=\frac{1}{4}\Sigma
\]

Vậy
\[
V_1\sim N_p\left(0,\frac{1}{4}\Sigma\right),
\qquad
V_2\sim N_p\left(0,\frac{1}{4}\Sigma\right).
\]

\end{answer}
\end{subquestion}

%--------2.7b--------
\begin{subquestion}
\begin{answer}
\textbf{Phân phối đồng thời của $V_1$ và $V_2$}

Ta tính hiệp phương sai
\[
\mathrm{Cov}(V_1,V_2)
\]

\[
=\frac{1}{16}(\Sigma-\Sigma-\Sigma+\Sigma)=0.
\]

Do đó ma trận hiệp phương sai của vector
\[
(V_1,V_2)'
\]

là
\[
\begin{pmatrix}
\frac{1}{4}\Sigma & 0 \\
0 & \frac{1}{4}\Sigma
\end{pmatrix}.
\]

Vì phân phối chuẩn nhiều chiều và hiệp phương sai bằng $0$ nên
\[
V_1 \perp V_2 .
\]

Do đó
\[
(V_1,V_2)' \sim
N_{2p}
\left(
\begin{pmatrix}
0\\
0
\end{pmatrix},
\begin{pmatrix}
\frac{1}{4}\Sigma & 0\\
0 & \frac{1}{4}\Sigma
\end{pmatrix}
\right).
\]
\end{answer}
\end{subquestion}

\end{question}